\documentclass[12pt]{article}
\newcommand{\bandlabel}{Grades 4--5}
\newcommand{\packetid}{CO-45-v1}
% Shared preamble for Week NN: Counting-out circles (final).
% Each band file defines \bandlabel and \packetid before % Shared preamble for Week NN: Counting-out circles (final).
% Each band file defines \bandlabel and \packetid before % Shared preamble for Week NN: Counting-out circles (final).
% Each band file defines \bandlabel and \packetid before \input{common}.
\usepackage[T1]{fontenc}
\usepackage[default]{sourcesanspro}
\usepackage[letterpaper,left=1.5cm,right=1.5cm,top=2.1cm,bottom=2.0cm,
  headheight=16pt,headsep=0.6cm,footskip=0.9cm]{geometry}
\usepackage{fancyhdr}
\usepackage{tikz}
\usetikzlibrary{calc}
\usepackage{xcolor}

\pagestyle{fancy}
\fancyhf{}
\renewcommand{\headrulewidth}{0pt}
\fancyhead[L]{\normalsize Week NN / Counting-out circles / \bandlabel}
\fancyfoot[L]{\small Bellingham Math Circle / Week NN / \packetid}
\fancyfoot[R]{\small\thepage}

\setlength{\parindent}{0pt}
\setlength{\parskip}{0.5\baselineskip}
\raggedright
\hyphenpenalty=10000
\exhyphenpenalty=10000

\definecolor{ringline}{gray}{0.72}
\definecolor{boxline}{gray}{0.55}

\newcommand{\prob}[1]{\textbf{Problem #1:}}

% \ringat{x}{y}{n}{R}{spot radius}{label gap}{label font}: draws inside the
% current tikzpicture. Spots are numbered 1..n clockwise with 1 at the top;
% the numbers sit just outside the spots so counters do not cover them.
\newcommand{\ringat}[7]{%
  \begin{scope}[shift={(#1,#2)}]
    \ifnum#3>1
      \draw[ringline, line width=0.8pt] (0,0) circle (#4);
    \fi
    \foreach \i in {1,...,#3} {
      \pgfmathsetmacro\ang{90-360*(\i-1)/#3}
      \filldraw[fill=white, draw=black, line width=1pt] (\ang:#4) circle (#5);
      \node[font=#7, inner sep=0pt] at (\ang:{#4+#5+#6}) {\i};
    }
  \end{scope}}

% Stand-alone ring picture: \ring{n}{R}{spot radius}{label gap}{label font}.
\newcommand{\ring}[5]{%
  \begin{tikzpicture}[baseline=(current bounding box.center)]
    \ringat{0}{0}{#1}{#2}{#3}{#4}{#5}
  \end{tikzpicture}}

% Working ring for counters (spots 2.6 cm across). Arg: n, R.
\newcommand{\bigring}[2]{\ring{#1}{#2}{1.3}{0.42}{\large\bfseries}}
\newcommand{\bigringat}[4]{\ringat{#1}{#2}{#3}{#4}{1.3}{0.42}{\large\bfseries}}

% A record table. #1 = first-column width (cm), #2 = cell width (cm),
% #3 = cell height (cm), #4 = comma list of column headings,
% #5 = row labels as {label1,label2,...} (first label goes with the headings row).
\newcommand{\rectable}[5]{%
  \begin{tikzpicture}[x=1cm,y=1cm]
    \def\fw{#1}\def\cw{#2}\def\ch{#3}
    \foreach \lab [count=\r from 0] in {#5} {
      \draw[line width=0.7pt] (0,-\r*\ch) rectangle (\fw,-\r*\ch-\ch);
      \node[anchor=west, inner sep=3pt, text width=\fw cm - 6pt, align=left,
            text height=1.6ex, text depth=0.4ex,
            font=\small] at (0,-\r*\ch-\ch/2) {\lab};
      \foreach \h [count=\c from 0] in {#4} {
        \draw[line width=0.7pt] (\fw+\c*\cw,-\r*\ch) rectangle (\fw+\c*\cw+\cw,-\r*\ch-\ch);
        \ifnum\r=0
          \node[font=\bfseries, text height=1.6ex, text depth=0.3ex] at (\fw+\c*\cw+\cw/2,-\ch/2) {\h};
        \fi
      }
    }
  \end{tikzpicture}}

% Answer box: #1 = width (cm), #2 = height (cm), #3 = optional label at top left.
\newcommand{\answerbox}[3]{%
  \begin{tikzpicture}
    \draw[boxline, line width=0.8pt, rounded corners=4pt] (0,0) rectangle (#1,-#2);
    \node[anchor=north west, inner sep=5pt, font=\small] at (0,0) {#3};
  \end{tikzpicture}}

% Full-page mat: a 16-spot ring with 2.6 cm spots (for 1-inch counters),
% about 18.5 cm across with its numbers. Used as a page of its own.
\newcommand{\matsixteen}{%
  \begin{tikzpicture}
    \ringat{0}{0}{16}{7.2}{1.3}{0.42}{\large\bfseries}
  \end{tikzpicture}}
.
\usepackage[T1]{fontenc}
\usepackage[default]{sourcesanspro}
\usepackage[letterpaper,left=1.5cm,right=1.5cm,top=2.1cm,bottom=2.0cm,
  headheight=16pt,headsep=0.6cm,footskip=0.9cm]{geometry}
\usepackage{fancyhdr}
\usepackage{tikz}
\usetikzlibrary{calc}
\usepackage{xcolor}

\pagestyle{fancy}
\fancyhf{}
\renewcommand{\headrulewidth}{0pt}
\fancyhead[L]{\normalsize Week NN / Counting-out circles / \bandlabel}
\fancyfoot[L]{\small Bellingham Math Circle / Week NN / \packetid}
\fancyfoot[R]{\small\thepage}

\setlength{\parindent}{0pt}
\setlength{\parskip}{0.5\baselineskip}
\raggedright
\hyphenpenalty=10000
\exhyphenpenalty=10000

\definecolor{ringline}{gray}{0.72}
\definecolor{boxline}{gray}{0.55}

\newcommand{\prob}[1]{\textbf{Problem #1:}}

% \ringat{x}{y}{n}{R}{spot radius}{label gap}{label font}: draws inside the
% current tikzpicture. Spots are numbered 1..n clockwise with 1 at the top;
% the numbers sit just outside the spots so counters do not cover them.
\newcommand{\ringat}[7]{%
  \begin{scope}[shift={(#1,#2)}]
    \ifnum#3>1
      \draw[ringline, line width=0.8pt] (0,0) circle (#4);
    \fi
    \foreach \i in {1,...,#3} {
      \pgfmathsetmacro\ang{90-360*(\i-1)/#3}
      \filldraw[fill=white, draw=black, line width=1pt] (\ang:#4) circle (#5);
      \node[font=#7, inner sep=0pt] at (\ang:{#4+#5+#6}) {\i};
    }
  \end{scope}}

% Stand-alone ring picture: \ring{n}{R}{spot radius}{label gap}{label font}.
\newcommand{\ring}[5]{%
  \begin{tikzpicture}[baseline=(current bounding box.center)]
    \ringat{0}{0}{#1}{#2}{#3}{#4}{#5}
  \end{tikzpicture}}

% Working ring for counters (spots 2.6 cm across). Arg: n, R.
\newcommand{\bigring}[2]{\ring{#1}{#2}{1.3}{0.42}{\large\bfseries}}
\newcommand{\bigringat}[4]{\ringat{#1}{#2}{#3}{#4}{1.3}{0.42}{\large\bfseries}}

% A record table. #1 = first-column width (cm), #2 = cell width (cm),
% #3 = cell height (cm), #4 = comma list of column headings,
% #5 = row labels as {label1,label2,...} (first label goes with the headings row).
\newcommand{\rectable}[5]{%
  \begin{tikzpicture}[x=1cm,y=1cm]
    \def\fw{#1}\def\cw{#2}\def\ch{#3}
    \foreach \lab [count=\r from 0] in {#5} {
      \draw[line width=0.7pt] (0,-\r*\ch) rectangle (\fw,-\r*\ch-\ch);
      \node[anchor=west, inner sep=3pt, text width=\fw cm - 6pt, align=left,
            text height=1.6ex, text depth=0.4ex,
            font=\small] at (0,-\r*\ch-\ch/2) {\lab};
      \foreach \h [count=\c from 0] in {#4} {
        \draw[line width=0.7pt] (\fw+\c*\cw,-\r*\ch) rectangle (\fw+\c*\cw+\cw,-\r*\ch-\ch);
        \ifnum\r=0
          \node[font=\bfseries, text height=1.6ex, text depth=0.3ex] at (\fw+\c*\cw+\cw/2,-\ch/2) {\h};
        \fi
      }
    }
  \end{tikzpicture}}

% Answer box: #1 = width (cm), #2 = height (cm), #3 = optional label at top left.
\newcommand{\answerbox}[3]{%
  \begin{tikzpicture}
    \draw[boxline, line width=0.8pt, rounded corners=4pt] (0,0) rectangle (#1,-#2);
    \node[anchor=north west, inner sep=5pt, font=\small] at (0,0) {#3};
  \end{tikzpicture}}

% Full-page mat: a 16-spot ring with 2.6 cm spots (for 1-inch counters),
% about 18.5 cm across with its numbers. Used as a page of its own.
\newcommand{\matsixteen}{%
  \begin{tikzpicture}
    \ringat{0}{0}{16}{7.2}{1.3}{0.42}{\large\bfseries}
  \end{tikzpicture}}
.
\usepackage[T1]{fontenc}
\usepackage[default]{sourcesanspro}
\usepackage[letterpaper,left=1.5cm,right=1.5cm,top=2.1cm,bottom=2.0cm,
  headheight=16pt,headsep=0.6cm,footskip=0.9cm]{geometry}
\usepackage{fancyhdr}
\usepackage{tikz}
\usetikzlibrary{calc}
\usepackage{xcolor}

\pagestyle{fancy}
\fancyhf{}
\renewcommand{\headrulewidth}{0pt}
\fancyhead[L]{\normalsize Week NN / Counting-out circles / \bandlabel}
\fancyfoot[L]{\small Bellingham Math Circle / Week NN / \packetid}
\fancyfoot[R]{\small\thepage}

\setlength{\parindent}{0pt}
\setlength{\parskip}{0.5\baselineskip}
\raggedright
\hyphenpenalty=10000
\exhyphenpenalty=10000

\definecolor{ringline}{gray}{0.72}
\definecolor{boxline}{gray}{0.55}

\newcommand{\prob}[1]{\textbf{Problem #1:}}

% \ringat{x}{y}{n}{R}{spot radius}{label gap}{label font}: draws inside the
% current tikzpicture. Spots are numbered 1..n clockwise with 1 at the top;
% the numbers sit just outside the spots so counters do not cover them.
\newcommand{\ringat}[7]{%
  \begin{scope}[shift={(#1,#2)}]
    \ifnum#3>1
      \draw[ringline, line width=0.8pt] (0,0) circle (#4);
    \fi
    \foreach \i in {1,...,#3} {
      \pgfmathsetmacro\ang{90-360*(\i-1)/#3}
      \filldraw[fill=white, draw=black, line width=1pt] (\ang:#4) circle (#5);
      \node[font=#7, inner sep=0pt] at (\ang:{#4+#5+#6}) {\i};
    }
  \end{scope}}

% Stand-alone ring picture: \ring{n}{R}{spot radius}{label gap}{label font}.
\newcommand{\ring}[5]{%
  \begin{tikzpicture}[baseline=(current bounding box.center)]
    \ringat{0}{0}{#1}{#2}{#3}{#4}{#5}
  \end{tikzpicture}}

% Working ring for counters (spots 2.6 cm across). Arg: n, R.
\newcommand{\bigring}[2]{\ring{#1}{#2}{1.3}{0.42}{\large\bfseries}}
\newcommand{\bigringat}[4]{\ringat{#1}{#2}{#3}{#4}{1.3}{0.42}{\large\bfseries}}

% A record table. #1 = first-column width (cm), #2 = cell width (cm),
% #3 = cell height (cm), #4 = comma list of column headings,
% #5 = row labels as {label1,label2,...} (first label goes with the headings row).
\newcommand{\rectable}[5]{%
  \begin{tikzpicture}[x=1cm,y=1cm]
    \def\fw{#1}\def\cw{#2}\def\ch{#3}
    \foreach \lab [count=\r from 0] in {#5} {
      \draw[line width=0.7pt] (0,-\r*\ch) rectangle (\fw,-\r*\ch-\ch);
      \node[anchor=west, inner sep=3pt, text width=\fw cm - 6pt, align=left,
            text height=1.6ex, text depth=0.4ex,
            font=\small] at (0,-\r*\ch-\ch/2) {\lab};
      \foreach \h [count=\c from 0] in {#4} {
        \draw[line width=0.7pt] (\fw+\c*\cw,-\r*\ch) rectangle (\fw+\c*\cw+\cw,-\r*\ch-\ch);
        \ifnum\r=0
          \node[font=\bfseries, text height=1.6ex, text depth=0.3ex] at (\fw+\c*\cw+\cw/2,-\ch/2) {\h};
        \fi
      }
    }
  \end{tikzpicture}}

% Answer box: #1 = width (cm), #2 = height (cm), #3 = optional label at top left.
\newcommand{\answerbox}[3]{%
  \begin{tikzpicture}
    \draw[boxline, line width=0.8pt, rounded corners=4pt] (0,0) rectangle (#1,-#2);
    \node[anchor=north west, inner sep=5pt, font=\small] at (0,0) {#3};
  \end{tikzpicture}}

% Full-page mat: a 16-spot ring with 2.6 cm spots (for 1-inch counters),
% about 18.5 cm across with its numbers. Used as a page of its own.
\newcommand{\matsixteen}{%
  \begin{tikzpicture}
    \ringat{0}{0}{16}{7.2}{1.3}{0.42}{\large\bfseries}
  \end{tikzpicture}}


% Problem 10 table: binary columns headed by place values.
\newcommand{\binarytable}{%
  \begin{tikzpicture}[x=1cm,y=1cm]
    \def\ch{1.0}\def\bw{1.0}
    % column x-positions: people [0,2.6], binary [2.6,8.6], last [9.0,12.4], binary [12.4,18.4]
    \foreach \x/\w/\lab in {0/2.6/people, 9.0/3.4/last one standing} {
      \draw[line width=0.7pt] (\x,0) rectangle (\x+\w,-2*\ch);
      \node[font=\small, text width=\w cm, align=center] at (\x+\w/2,-\ch) {\lab};
      \foreach \r in {2,...,9} {
        \draw[line width=0.7pt] (\x,-\r*\ch) rectangle (\x+\w,-\r*\ch-\ch);
      }
    }
    \foreach \x in {2.6, 12.4} {
      \draw[line width=0.7pt] (\x,0) rectangle (\x+6*\bw,-\ch);
      \node[font=\small] at (\x+3*\bw,-\ch/2) {in binary};
      \foreach \pv [count=\c from 0] in {32,16,8,4,2,1} {
        \draw[line width=0.7pt] (\x+\c*\bw,-\ch) rectangle (\x+\c*\bw+\bw,-2*\ch);
        \node[font=\footnotesize\bfseries] at (\x+\c*\bw+\bw/2,-1.5*\ch) {\pv};
        \foreach \r in {2,...,9} {
          \draw[line width=0.7pt] (\x+\c*\bw,-\r*\ch) rectangle (\x+\c*\bw+\bw,-\r*\ch-\ch);
        }
      }
    }
    % worked row: 12 = 1100, last one standing 9 = 1001
    \node[font=\large] at (1.3,-2.5*\ch) {12};
    \foreach \d [count=\c from 2] in {1,1,0,0} {
      \node[font=\large] at (2.6+\c*\bw+\bw/2,-2.5*\ch) {\d};
    }
    \node[font=\large] at (10.7,-2.5*\ch) {9};
    \foreach \d [count=\c from 2] in {1,0,0,1} {
      \node[font=\large] at (12.4+\c*\bw+\bw/2,-2.5*\ch) {\d};
    }
  \end{tikzpicture}}

\begin{document}

Each counter is a person. To play a circle of 10 people, put counters on spots 1 to 10 of a mat and leave the other spots empty. Start at 1 and go round and round the circle. The people standing take turns: 1 stays standing, the next one sits down, the next one stays standing, the next one sits down, and so on. When you come round past 1 again, keep taking turns from where you left off. Flip a counter over when that person sits down. Skip empty spots and people who are sitting. Stop when only one person is standing.

\prob{1} Play circles of 1 to 16 people on the mat with 16 spots. Write the last one standing for each circle.

\noindent\rectable{3.6}{1.7}{1.05}{1,2,3,4,5,6,7,8}{people,last one standing}

\vspace{0.3cm}
\noindent\rectable{3.6}{1.7}{1.05}{9,10,11,12,13,14,15,16}{people,last one standing}

\vspace{0.4cm}
\prob{2} Predict the last one standing in circles of 20, 25 and 32 people. Then play them on the mat with 32 spots to check.

\noindent\rectable{3.6}{2.4}{1.05}{20,25,32}{people,prediction,last one standing}

\vspace{0.4cm}
\prob{3} Find every circle up to 100 people where number 1 is the last one standing. Explain why number 1 is last in all of these circles.

\noindent\begin{tikzpicture}
  \draw[boxline, line width=0.8pt, rounded corners=4pt] (0,0) rectangle (18.5,-7.6);
  \ringat{4.2}{-3.8}{16}{2.75}{0.33}{0.27}{\small}
\end{tikzpicture}

\newpage
\prob{4} Who is the last one standing in a circle of 100 people? In a circle of 1000 people? Explain why your answers are right.

\noindent\answerbox{18.5}{9.5}{}

\vspace{0.5cm}
\prob{5} Number 21 is the last one standing in a circle of 26 people. Explain how this tells you who is last in a circle of 52 people, and who is last in a circle of 53 people.

\noindent\answerbox{18.5}{9.5}{}

\newpage
\prob{6} Find every circle up to 100 people where number 23 is the last one standing. Do the same for number 22. If there are none, explain why.

\noindent\answerbox{18.5}{4.2}{Number 23:}

\vspace{0.3cm}
\noindent\answerbox{18.5}{5.2}{Number 22:}

\vspace{0.5cm}
\prob{7} Now every third person sits down: 1 and 2 stay standing, the next one sits down, the next two stay standing, the next one sits down, and so on. Play circles of 1 to 12 people and write the last one standing for each.

\noindent\rectable{3.6}{2.2}{1.15}{1,2,3,4,5,6}{people,last one standing}

\vspace{0.3cm}
\noindent\rectable{3.6}{2.2}{1.15}{7,8,9,10,11,12}{people,last one standing}

\newpage
\prob{8} With every third person sitting down, predict the last one standing in circles of 13 to 24 people. Then play circles of 14 and 21 to check.

\noindent\rectable{3.6}{2.2}{1.1}{13,14,15,16,17,18}{people,prediction,last one standing}

\vspace{0.3cm}
\noindent\rectable{3.6}{2.2}{1.1}{19,20,21,22,23,24}{people,prediction,last one standing}

\vspace{0.5cm}
\prob{9} In a circle of 41 people, every third person sits down. Where should you stand to be the last one standing?

\noindent\answerbox{18.5}{12.0}{}

\newpage
\prob{10} With every second person sitting down, the last one standing in a circle of 12 people is 9. In binary, 12 is written 1100, because \mbox{12 = 8 + 4}, and 9 is written 1001, because \mbox{9 = 8 + 1}. Find a rule that turns the number of people, written in binary, into the last one standing, written in binary. Explain why your rule works.

\vspace{0.2cm}
\noindent\binarytable

\vspace{0.5cm}
\noindent\answerbox{18.5}{6.5}{}

\newpage
\noindent Mat with 16 spots.\par
\vspace*{\fill}
\begin{center}
\matsixteen
\end{center}
\vspace*{\fill}

\newpage
\noindent% generated by make_mat32.py; do not edit by hand
\begin{tikzpicture}[x=1cm,y=1cm]
\useasboundingbox (0,0.30) rectangle (18.59,23.54);
\draw[ringline, line width=0.8pt] (21.090,2.250) -- (5.400,2.250) arc[start angle=270, end angle=180, radius=3.00] -- (2.400,18.590) arc[start angle=180, end angle=90, radius=3.00] -- (21.090,21.590);
\filldraw[fill=white, draw=black, line width=1pt] (17.990,2.250) circle (1.30);
\node[font=\large\bfseries, inner sep=0pt] at (17.990,0.330) {22};
\filldraw[fill=white, draw=black, line width=1pt] (14.794,2.250) circle (1.30);
\node[font=\large\bfseries, inner sep=0pt] at (14.794,0.330) {23};
\filldraw[fill=white, draw=black, line width=1pt] (11.597,2.250) circle (1.30);
\node[font=\large\bfseries, inner sep=0pt] at (11.597,0.330) {24};
\filldraw[fill=white, draw=black, line width=1pt] (8.401,2.250) circle (1.30);
\node[font=\large\bfseries, inner sep=0pt] at (8.401,0.330) {25};
\filldraw[fill=white, draw=black, line width=1pt] (5.205,2.256) circle (1.30);
\node[font=\large\bfseries, inner sep=0pt] at (5.080,0.340) {26};
\filldraw[fill=white, draw=black, line width=1pt] (2.686,3.971) circle (1.30);
\node[font=\large\bfseries, inner sep=0pt] at (0.949,3.153) {27};
\filldraw[fill=white, draw=black, line width=1pt] (2.400,7.126) circle (1.30);
\node[font=\large\bfseries, inner sep=0pt] at (0.480,7.126) {28};
\filldraw[fill=white, draw=black, line width=1pt] (2.400,10.322) circle (1.30);
\node[font=\large\bfseries, inner sep=0pt] at (0.480,10.322) {29};
\filldraw[fill=white, draw=black, line width=1pt] (2.400,13.518) circle (1.30);
\node[font=\large\bfseries, inner sep=0pt] at (0.480,13.518) {30};
\filldraw[fill=white, draw=black, line width=1pt] (2.400,16.714) circle (1.30);
\node[font=\large\bfseries, inner sep=0pt] at (0.480,16.714) {31};
\filldraw[fill=white, draw=black, line width=1pt] (2.686,19.869) circle (1.30);
\node[font=\large\bfseries, inner sep=0pt] at (0.949,20.687) {32};
\filldraw[fill=white, draw=black, line width=1pt] (5.205,21.584) circle (1.30);
\node[font=\large\bfseries, inner sep=0pt] at (5.080,23.500) {1};
\filldraw[fill=white, draw=black, line width=1pt] (8.401,21.590) circle (1.30);
\node[font=\large\bfseries, inner sep=0pt] at (8.401,23.510) {2};
\filldraw[fill=white, draw=black, line width=1pt] (11.597,21.590) circle (1.30);
\node[font=\large\bfseries, inner sep=0pt] at (11.597,23.510) {3};
\filldraw[fill=white, draw=black, line width=1pt] (14.794,21.590) circle (1.30);
\node[font=\large\bfseries, inner sep=0pt] at (14.794,23.510) {4};
\filldraw[fill=white, draw=black, line width=1pt] (17.990,21.590) circle (1.30);
\node[font=\large\bfseries, inner sep=0pt] at (17.990,23.510) {5};
\node[align=center, text width=7cm] at (9.295,11.920) {Mat with 32 spots, left half. Lay the right half beside it.};
\end{tikzpicture}


\newpage
\noindent% generated by make_mat32.py; do not edit by hand
\begin{tikzpicture}[x=1cm,y=1cm]
\useasboundingbox (0,0.30) rectangle (18.59,23.54);
\draw[ringline, line width=0.8pt] (-2.500,2.250) -- (13.190,2.250) arc[start angle=270, end angle=360, radius=3.00] -- (16.190,18.590) arc[start angle=0, end angle=90, radius=3.00] -- (-2.500,21.590);
\filldraw[fill=white, draw=black, line width=1pt] (0.600,2.250) circle (1.30);
\node[font=\large\bfseries, inner sep=0pt] at (0.600,0.330) {21};
\filldraw[fill=white, draw=black, line width=1pt] (3.796,2.250) circle (1.30);
\node[font=\large\bfseries, inner sep=0pt] at (3.796,0.330) {20};
\filldraw[fill=white, draw=black, line width=1pt] (6.993,2.250) circle (1.30);
\node[font=\large\bfseries, inner sep=0pt] at (6.993,0.330) {19};
\filldraw[fill=white, draw=black, line width=1pt] (10.189,2.250) circle (1.30);
\node[font=\large\bfseries, inner sep=0pt] at (10.189,0.330) {18};
\filldraw[fill=white, draw=black, line width=1pt] (13.385,2.256) circle (1.30);
\node[font=\large\bfseries, inner sep=0pt] at (13.510,0.340) {17};
\filldraw[fill=white, draw=black, line width=1pt] (15.904,3.971) circle (1.30);
\node[font=\large\bfseries, inner sep=0pt] at (17.641,3.153) {16};
\filldraw[fill=white, draw=black, line width=1pt] (16.190,7.126) circle (1.30);
\node[font=\large\bfseries, inner sep=0pt] at (18.110,7.126) {15};
\filldraw[fill=white, draw=black, line width=1pt] (16.190,10.322) circle (1.30);
\node[font=\large\bfseries, inner sep=0pt] at (18.110,10.322) {14};
\filldraw[fill=white, draw=black, line width=1pt] (16.190,13.518) circle (1.30);
\node[font=\large\bfseries, inner sep=0pt] at (18.110,13.518) {13};
\filldraw[fill=white, draw=black, line width=1pt] (16.190,16.714) circle (1.30);
\node[font=\large\bfseries, inner sep=0pt] at (18.110,16.714) {12};
\filldraw[fill=white, draw=black, line width=1pt] (15.904,19.869) circle (1.30);
\node[font=\large\bfseries, inner sep=0pt] at (17.641,20.687) {11};
\filldraw[fill=white, draw=black, line width=1pt] (13.385,21.584) circle (1.30);
\node[font=\large\bfseries, inner sep=0pt] at (13.510,23.500) {10};
\filldraw[fill=white, draw=black, line width=1pt] (10.189,21.590) circle (1.30);
\node[font=\large\bfseries, inner sep=0pt] at (10.189,23.510) {9};
\filldraw[fill=white, draw=black, line width=1pt] (6.993,21.590) circle (1.30);
\node[font=\large\bfseries, inner sep=0pt] at (6.993,23.510) {8};
\filldraw[fill=white, draw=black, line width=1pt] (3.796,21.590) circle (1.30);
\node[font=\large\bfseries, inner sep=0pt] at (3.796,23.510) {7};
\filldraw[fill=white, draw=black, line width=1pt] (0.600,21.590) circle (1.30);
\node[font=\large\bfseries, inner sep=0pt] at (0.600,23.510) {6};
\node[align=center, text width=7cm] at (9.295,11.920) {Mat with 32 spots, right half.};
\end{tikzpicture}


\end{document}
